% Options for packages loaded elsewhere
% Options for packages loaded elsewhere
\PassOptionsToPackage{unicode}{hyperref}
\PassOptionsToPackage{hyphens}{url}
\PassOptionsToPackage{dvipsnames,svgnames,x11names}{xcolor}
%
\documentclass[
  letterpaper,
  DIV=11,
  numbers=noendperiod]{scrartcl}
\usepackage{xcolor}
\usepackage{amsmath,amssymb}
\setcounter{secnumdepth}{5}
\usepackage{iftex}
\ifPDFTeX
  \usepackage[T1]{fontenc}
  \usepackage[utf8]{inputenc}
  \usepackage{textcomp} % provide euro and other symbols
\else % if luatex or xetex
  \usepackage{unicode-math} % this also loads fontspec
  \defaultfontfeatures{Scale=MatchLowercase}
  \defaultfontfeatures[\rmfamily]{Ligatures=TeX,Scale=1}
\fi
\usepackage{lmodern}
\ifPDFTeX\else
  % xetex/luatex font selection
\fi
% Use upquote if available, for straight quotes in verbatim environments
\IfFileExists{upquote.sty}{\usepackage{upquote}}{}
\IfFileExists{microtype.sty}{% use microtype if available
  \usepackage[]{microtype}
  \UseMicrotypeSet[protrusion]{basicmath} % disable protrusion for tt fonts
}{}
\makeatletter
\@ifundefined{KOMAClassName}{% if non-KOMA class
  \IfFileExists{parskip.sty}{%
    \usepackage{parskip}
  }{% else
    \setlength{\parindent}{0pt}
    \setlength{\parskip}{6pt plus 2pt minus 1pt}}
}{% if KOMA class
  \KOMAoptions{parskip=half}}
\makeatother
% Make \paragraph and \subparagraph free-standing
\makeatletter
\ifx\paragraph\undefined\else
  \let\oldparagraph\paragraph
  \renewcommand{\paragraph}{
    \@ifstar
      \xxxParagraphStar
      \xxxParagraphNoStar
  }
  \newcommand{\xxxParagraphStar}[1]{\oldparagraph*{#1}\mbox{}}
  \newcommand{\xxxParagraphNoStar}[1]{\oldparagraph{#1}\mbox{}}
\fi
\ifx\subparagraph\undefined\else
  \let\oldsubparagraph\subparagraph
  \renewcommand{\subparagraph}{
    \@ifstar
      \xxxSubParagraphStar
      \xxxSubParagraphNoStar
  }
  \newcommand{\xxxSubParagraphStar}[1]{\oldsubparagraph*{#1}\mbox{}}
  \newcommand{\xxxSubParagraphNoStar}[1]{\oldsubparagraph{#1}\mbox{}}
\fi
\makeatother


\usepackage{longtable,booktabs,array}
\usepackage{calc} % for calculating minipage widths
% Correct order of tables after \paragraph or \subparagraph
\usepackage{etoolbox}
\makeatletter
\patchcmd\longtable{\par}{\if@noskipsec\mbox{}\fi\par}{}{}
\makeatother
% Allow footnotes in longtable head/foot
\IfFileExists{footnotehyper.sty}{\usepackage{footnotehyper}}{\usepackage{footnote}}
\makesavenoteenv{longtable}
\usepackage{graphicx}
\makeatletter
\newsavebox\pandoc@box
\newcommand*\pandocbounded[1]{% scales image to fit in text height/width
  \sbox\pandoc@box{#1}%
  \Gscale@div\@tempa{\textheight}{\dimexpr\ht\pandoc@box+\dp\pandoc@box\relax}%
  \Gscale@div\@tempb{\linewidth}{\wd\pandoc@box}%
  \ifdim\@tempb\p@<\@tempa\p@\let\@tempa\@tempb\fi% select the smaller of both
  \ifdim\@tempa\p@<\p@\scalebox{\@tempa}{\usebox\pandoc@box}%
  \else\usebox{\pandoc@box}%
  \fi%
}
% Set default figure placement to htbp
\def\fps@figure{htbp}
\makeatother





\setlength{\emergencystretch}{3em} % prevent overfull lines

\providecommand{\tightlist}{%
  \setlength{\itemsep}{0pt}\setlength{\parskip}{0pt}}



 


\KOMAoption{captions}{tableheading}
\makeatletter
\@ifpackageloaded{caption}{}{\usepackage{caption}}
\AtBeginDocument{%
\ifdefined\contentsname
  \renewcommand*\contentsname{Table of contents}
\else
  \newcommand\contentsname{Table of contents}
\fi
\ifdefined\listfigurename
  \renewcommand*\listfigurename{List of Figures}
\else
  \newcommand\listfigurename{List of Figures}
\fi
\ifdefined\listtablename
  \renewcommand*\listtablename{List of Tables}
\else
  \newcommand\listtablename{List of Tables}
\fi
\ifdefined\figurename
  \renewcommand*\figurename{Figure}
\else
  \newcommand\figurename{Figure}
\fi
\ifdefined\tablename
  \renewcommand*\tablename{Table}
\else
  \newcommand\tablename{Table}
\fi
}
\@ifpackageloaded{float}{}{\usepackage{float}}
\floatstyle{ruled}
\@ifundefined{c@chapter}{\newfloat{codelisting}{h}{lop}}{\newfloat{codelisting}{h}{lop}[chapter]}
\floatname{codelisting}{Listing}
\newcommand*\listoflistings{\listof{codelisting}{List of Listings}}
\makeatother
\makeatletter
\makeatother
\makeatletter
\@ifpackageloaded{caption}{}{\usepackage{caption}}
\@ifpackageloaded{subcaption}{}{\usepackage{subcaption}}
\makeatother
\usepackage{bookmark}
\IfFileExists{xurl.sty}{\usepackage{xurl}}{} % add URL line breaks if available
\urlstyle{same}
\hypersetup{
  pdftitle={Lab 7: Curvy data},
  pdfauthor={Jeremy Hemberger and Brian Aukema},
  colorlinks=true,
  linkcolor={blue},
  filecolor={Maroon},
  citecolor={Blue},
  urlcolor={Blue},
  pdfcreator={LaTeX via pandoc}}


\title{Lab 7: Curvy data}
\usepackage{etoolbox}
\makeatletter
\providecommand{\subtitle}[1]{% add subtitle to \maketitle
  \apptocmd{\@title}{\par {\large #1 \par}}{}{}
}
\makeatother
\subtitle{Entomology 5126 - Spring 2026}
\author{Jeremy Hemberger and Brian Aukema}
\date{March 6, 2026}
\begin{document}
\maketitle

\renewcommand*\contentsname{Table of contents}
{
\hypersetup{linkcolor=}
\setcounter{tocdepth}{3}
\tableofcontents
}

\begin{center}\rule{0.5\linewidth}{0.5pt}\end{center}

\section{Introduction}\label{introduction}

Today, we'll work on regression techniques for curvy data. Sometimes,
when you first graph the data, you can see instantly that there is an
asymptote, or a gentle curve, or a large squiggle, or\ldots. whatever it
is, a straight line will not be a good model. And let's face it, there
are very few patterns in ecology that are straight lines! Today, we will
focus on two simple techniques:

\begin{itemize}
\item
  Transforming an independent (\(x\)) variable. We fit a straight line
  to data, that, when back-transformed, becomes curved!
\item
  Polynomial regression. This technique is a special type of multiple
  regression. Of course, you are good at multiple regression now.
\end{itemize}

\section{Fitting curves}\label{fitting-curves}

\subsection{\texorpdfstring{Transforming the \(x\)
variable}{Transforming the x variable}}\label{transforming-the-x-variable}

Decay functions and asymptotes are frequently seen in dose-response
curves where increasing doses have decreasing biological effects, or
perhaps an effect decays through time. One of the easiest things to do
is transform the \(x\) variable, fit a regular regression line, and then
back-transform.

Reciprocal \(x\) transformation

\pandocbounded{\includegraphics[keepaspectratio]{../content/slides/slide_media/w7/reciprocal-curve.png}}

and here

More text

\(y_i = \beta_0 + \beta_1 (1/x_i) + \epsilon_i\)
\(y_i = \beta_0 + \beta_1\)log\((x_i) + \epsilon_i\)

\subsection{Polynomial regression}\label{polynomial-regression}

The basic formula for polynomial regression is
\[y_i = \beta_0 + \beta_1 x + \beta_2 x^{2} + \ldots + \beta_k x^{k} + \epsilon_i\]
The \(\beta_0\) is the intercept, or where the line crosses the
\(y\)-axis. The \(\beta_1\) is the slope of the first covariate, the
``linear'' term, or the increase in \(y\) per unit change in \(x\)
\emph{holding all other covariates constant}. The \(\beta_2\) is the
slope of the second covariate, the ``quadratic'' term. The
\(\epsilon_i\) refers to the ``errors'' about the line, and we assume
that they are normally distributed with a constant variance (i.e.,
\(\epsilon_i \sim N(0, \sigma^2)\)).\\
An important note when reporting a polynomial regression line: by
convention, we always report ``lower order'' terms, even if they are not
significant. So, if you have a quadratic equation, for example, where
only the quadratic (i.e., squared) covariate is significant, we would
still include the intercept and the linear term when reporting the
equation of the line, even if they are not significantly different from
zero.

\subsection{Checking assumptions}\label{checking-assumptions}

Both techniques have the same linear model structure. As such, you are
familiar with the assumptions:

\begin{enumerate}
\def\labelenumi{\arabic{enumi}.}
\item
  The model is correct
\item
  The observations are independent
\item
  The variances are equal
\item
  The errors are normally distributed
\end{enumerate}

As you might expect, you can check the assumptions of equal variances
and normally distributed errors in your residual plot, as well as
detecting if your model is perhaps misspecified.

\subsection{\texorpdfstring{Transforming a covariate in
\texttt{R}}{Transforming a covariate in R}}\label{transforming-a-covariate-in-r}

Both of these techniques involve transforming a covariate. For example,
in polynomial regression, you need to square a covariate to get the
quadratic term. When doing this in \texttt{R}, you need to use the
\texttt{identity} specifier, \texttt{I()}. That's just a fancy way of
telling the program to transform and protect the covariate. For example:

\begin{verbatim}
> fm1 <- lm(y ~ log(x), ...) # standard function, no identity protection needed
> fm2 <- lm(y ~ I(1/x), ...) # add data argument as before
> fm3 <- lm(y ~ x + I(x^2), ...) # quadratic function
\end{verbatim}

\subsection{Other techniques (FYI)}\label{other-techniques-fyi}

There are \emph{many} other techniques for fitting curves: loess fitting
(also known as \emph{locally weighted polynomial regression} or
\emph{locally weighted regression scatterplot smoothing}), smoothing
splines, and \emph{nonlinear} models.

For example, the asymptotic decay function below

could also be fit using a nonlinear model like this:

\(y(x) = \phi_1 + (\phi_2-\phi_1)\)exp\([-\)exp\((\phi_3)x]\)

I am not going to cover nonlinear methods in this course. More often
than not, the practitioner fits a nonlinear model because it looks good
and the software offered the option (maybe the button was shiny?).
Keeping with a theme in the course, \textbf{DON'T do something unless
you can explain what you did and why you did it}. From the example
above, one may arrive at a nice line, but if you ask what the estimate
of \(\phi_3\) represents (is it an intercept? a slope?), the
practitioner may not have a clue. That's not the way you want to conduct
your analyses. (In the above example, it's actually the logarithm of the
rate constant. A half-life could be estimated \(t_{0.5} =\) log
\(2/\)exp\((\phi_3)\)).

The nice thing about using a ``regular'' regression line (even with a
transformed \(x\) variable) is that the form of the equation is still
very straightforward.

\section{Today's data}\label{todays-data}

Today's data set consists of the normal mean minimum temperatures in
January (based on 30 years of data from 1931-1960) for 56 cities across
the continental United States. It can be found in Hand, D.J, et
al.~(1984) \emph{A Handbook of Small Data Sets}, London: Chapman Hall,
208-210.

Variable Description

\begin{center}\rule{0.5\linewidth}{0.5pt}\end{center}

City Yeah, that State State postal abbreviation JanTemp Average minimum
January temperature in degrees Farenheit Lat Latitude; degrees north of
the equator Lon Longitude; degrees west of the prime meridian

\section{Today's Assignment}\label{todays-assignment}

\begin{enumerate}
\def\labelenumi{\arabic{enumi}.}
\item
  First, let's analyze some mock data. To get it into \texttt{R}, you
  can put it into an Excel spreadsheet for export, or enter it directly.

\begin{verbatim}
> x <- 1:10 # or, x <- c(1,2,3,4,5,6,7,8,9,10) is same thing
> y <- c(10,5,3.5,1.5,1,1,0.9,0.9,0.5,0.4)
> mydata <- data.frame(x=x, y=y) # make dataframe with these
>                                #  numbers, use same titles
> mydata
...
\end{verbatim}

  Because it is second nature, I do not need to remind you to plot the
  data (\(y\) vs.~\(x\)) and think about what sort of model to fit.
  Because the only regression line you have learned how to fit up to
  this point is a straight line, go ahead and do so now (use the
  \texttt{lm()} function as before).

  Report your results. Do \emph{not} simply reproduce the
  \texttt{summary()} table: tell me whether the model is significant
  (\(F\), \(P\)-values?) and comment on the \(R^2\) value. If you had
  not initially graphed the data, would you be satisfied with this
  result?
\item
  Examine and comment on the residual plot, particularly on any
  analytical assumptions that might be violated.
\item
  : This is where you exercise the creative, right half of your brain
  after a perhaps stressful midterm. Connect the dots in your residual
  plot and draw something that will make Audrey and I smile when marking
  the labs. You may add a few sentences of explanation if helpful.
\item
  Now fit a reciprocal transformation of the \(x\) variable (a decay
  fit). Examine and comment on the residual plot from the reciprocal
  transformation and the analytical results. Is this any better than the
  straight line fit?
\item
  Let's try a different sort of nonlinear curve. Create a new response
  variable by subtracting all the values from 10 as follows:

\begin{verbatim}
> mydata$new_y <- 10 - mydata$y
\end{verbatim}

  Plot the data, and fit a straight regression line (you do not need to
  include this in your assignment). You will notice that the residual
  plot indicates a poor model fit.\\
  Fit a logarithmic transformation of \(x\) instead. Does the \(R^2\)
  value improve? Have we captured more of the variation than using a
  simple straight line?

  (Actually, the residual plot still suggests that we have not fully
  captured the curve, but, for teaching and learning purposes, set this
  fact aside for now.)
\item
  Ok, as we have covered, one of the nice things about the previous
  transformations is that we are still using a simple, intuitive
  regression line

  \(y_i \sim \beta_0 + \beta_1x_i + \epsilon_i\), where
  \(\epsilon_i \sim N(0,\sigma^2)\)

  The covariate may be transformed, but the form of the equation is
  still pretty straightforward.

  That doesn't, however, prevent some of us from being easily confused:
  what \emph{does} the line of best fit look like, and how do we plot
  it?

  You have an \(x\) data range of 1 to 10. Let's say we want the
  predicted value (\(\hat{y}\)) at \(x=5.5\). The natural log of 5.5 is
  1.704748. If I multiply that number by the slope, and then add the
  intercept, I get my new \(y\) value to graph against \(x=5.5\).
  Simple.

  But, if I was really graphing this, I wouldn't want to do all these
  calculations by hand. There \emph{has} to be an easier way to do this,
  right?

  Of course there is. In \texttt{R}, you can extract the fitted value
  for each value of \(x\) by using \texttt{fitted.values(model)}. For
  example,

\begin{verbatim}
> plot(mydata$new_y ~ mydata$x) # plot data
> fitted.values(fm2) # model with log(x) from previous question
> lines(fitted.values(fm2)) # overplot y-hat
\end{verbatim}

  That method is useful for predicting values at all of the existing
  values of your \(x\) variable in the dataset. Right now, those values
  are 1 to 10. If I was going to add a line of best fit in my graphing
  software (SigmaPlot, Excel, whatever), I might specify a range from 0
  to 10, by increments of 0.1 (or even finer, so it is a nice smooth
  curve). A more complex method, but useful for getting your estimated
  \(\hat{y}\) for \emph{any} value of \(x\) (instead of at only the
  existing \(x\) values in your data set), is to use the
  \texttt{predict()} command. First, create a data frame of what \(x\)
  values you want to use for predictions. It can be as long or as short
  as you want (continued next page).

\begin{verbatim}
> new.data <- data.frame(x=1:100/10) # what does this do?
> new.data                           # oh! cool!
\end{verbatim}

  It is \textbf{critical} to name the \(x\) variable the \emph{same name
  as the} \(x\) covariate in your model. Otherwise, the prediction
  command, referring to the model fit, won't have a bloomin' clue what
  covariate you are talking about. So, if you fit
  \texttt{y}\(\sim\)\texttt{temperature} in your model formula, you
  should be naming the column in the data frame above
  \texttt{temperature} and not \texttt{x}, ok? To continue the example,

\begin{verbatim}
> mypredictions <- predict(fm2, new.data, se.fit=T)
\end{verbatim}

  This will work for \emph{any} regression line (as long as you specify
  the right covariates). If you are using multiple regression, specify
  multiple new columns in your data frame, corresponding to the desired
  names of your covariates.

  The resulting object, \texttt{mypredictions}, yields four different
  internal objects of varying lengths as part of the output. For
  example, the predictions store information about how many degrees of
  freedom were in the error term of the original model and the residual
  scale. We are usually most interested in the \texttt{fit}, and the
  standard errors of the predicted \(y\)-values, the \texttt{se.fit}.

  Let's place them in columns alongside the new \(x\) values:

\begin{verbatim}
> new.data$y.hats <- mypredictions$fit
> new.data$y.hat.SEs <- mypredictions$se.fit
> new.data
\end{verbatim}

  Provide the \(\hat{y}\) and the associated prediction error for a new
  \(x\) value at \(x=3.5\). (Please use the line from question 5, the
  log(\(x\)) transform, and not the reciprocal transform from question
  4).
\item
  Finally, we get to the real data for today; we turn our attention to
  polynomial fitting. Read in the data, and do a scatterplot matrix.\\
  We are finally working with spatial data! First, what strikes you as
  strange in the scatterplot of \texttt{Lat} on the \(y\) axis and
  \texttt{Long} on the \(x\) axis? Why does this occur? (Not required,
  but can you find a quick fix?)
\item
  Find and report a good spatial regression equation explaining the
  minimum January temperatures. Explain, from both ecological and
  statistical points of view, why you chose the final terms that you
  did.
\end{enumerate}




\end{document}
